\section{Preliminaries}

In this section, we provide the necessary background and introduce some notations that we shall use in this paper.

Throughout, we fix a field $\F$ of arbitrary characteristic.

\subsection{Symmetric groups}

Let $X$ be a finite set. The symmetric group $\sym{X}$ on $X$ is the group of bijections from $X$ to $X$ under composition of functions. By convention, $\sym{\varnothing}$ is the trivial group. When $Y$ is a non-empty subset of $X$, we view $\sym{Y}$ as a subgroup of $\sym{X}$ by identifying an element of $\sym{Y}$ with its extension that sends $x$ to $x$ for all $x \in X \setminus Y$.

Let $X \subseteq \mathbb{Z}^+$, and $k \in \mathbb{Z}^+$. Define the $k$-translated subset $\tl{X}{k}$ of $\mathbb{Z}^+$ by
\begin{gather*}\tl{X}{k} = \{ x+ k\colon x \in X \}. \end{gather*}
Let $\tlsym{X}{k}=\big\{\tl{f}{k}\colon f\in\sym{X}\big\}$ where $\tl{f}{k}\colon \tl{X}{k}\to\tl{X}{k}$ is defined by
\begin{gather*}
\tl{f}{k}(x+k) = f(x) + k
\end{gather*}
for all $x \in X$. Clearly, $f \mapsto \tl{f}{k}$ is a group isomorphism from $\sym{X}$ to $\tlsym{X}{k}$. For any $S \subseteq \sym{X}$, write $\tl{S}{k}$ for $\{ \tl{s}{k}\colon s \in S\}$.

For $n \in \mathbb{Z}^+$, we write $\sym{n}$ for $\sym{\{1,2,\dotsc, n\}}$, the usual symmetric group on $n$ letters.

\subsection{Compositions and Young subgroups}

A composition is a finite sequence of positive integers. Let $\lambda = (\lambda_1,\lambda_2,\dotsc, \lambda_r)$ be a composition. Then $\ell(\lambda) = r$. If $\sum\limits_{i=1}^{\ell(\lambda)} \lambda_i = n$, we say that $\lambda$ is a composition of $n$, and write $|\lambda| = n$. If in addition $\lambda_1 \geq \lambda_2 \geq \dotsb \geq \lambda_{\ell(\lambda)}$, then we call $\lambda$ a partition of $n$. The unique composition (or partition) of 0 is denoted by $\varnothing$, and $\ell(\varnothing)=0$.

To a composition $\lambda$ of $n$, we associate the Young subgroup $\sym{\lambda}$ of $\sym{n}$, where{\samepage
\begin{gather*}
\sym{\lambda} = \sym{\lambda_1} \tlsym{\lambda_2}{s(\lambda)_1} \dotsm \tlsym{\lambda_{\ell(\lambda)}}{s(\lambda)_{\ell(\lambda)-1}} \cong \sym{\lambda_1} \times \sym{\lambda_2} \times \dotsb \times \sym{\lambda_{\ell(\lambda)}}
\end{gather*}
and $s(\lambda)_j = \sum\limits_{i=1}^j \lambda_i$ for all $1 \leq j \leq \ell(\lambda)-1$.}

Let $(\alpha|\beta)$ be a bicomposition of $n$, i.e., $\alpha$ and $\beta$ are compositions and $|\alpha| + |\beta| = n$. One may view $(\alpha|\beta)$ as a composition of $n$ by concatenating $\alpha$ with $\beta$, and so we have the associated Young subgroup
\begin{gather*}
\sym{\alpha|\beta} = \sym{\alpha} \tlsym{\beta}{|\alpha|} \cong \sym{\alpha} \times \sym{\beta}.\end{gather*}

\subsection{Young diagrams}

For a partition $\lambda$, its Young diagram $[\lambda]$ is defined as
\begin{gather*}
[\lambda] = \big\{ (i,j) \in \mathbb{Z}^2 \colon 1\leq i \leq \ell(\lambda),\ 1 \leq j \leq \lambda_i \big\}.
\end{gather*}

The conjugate partition of $\lambda$, denoted $\lambda'$, is the partition whose Young diagram $[\lambda'] = \{ (i,j)\colon (j,i) \in [\lambda] \}$.

\subsection{Tableaux} \label{SS:Tableaux}

Let $\lambda$ be a partition of $n$. Define formally a $\lambda$-tableau $\t$ as a bijective function $\t\colon [\lambda] \to \{1,2,\dotsc,n \}$. We usually view this as a labelling of the elements of $[\lambda]$ by numbers $1,2,\dotsc,n$, such that each number appears exactly once. Denote the set of $\lambda$-tableaux by $\TT(\lambda)$.

Through a fixed $\lambda$-tableau $\t$ we obtain a group isomorphism $\sym{n} \cong \sym{[\lambda]}$ where we identify $\sigma\in\sym{n}$ with $\t^{-1} \circ \sigma \circ \t\in \sym{[\lambda]}$. Let $R_{[\lambda]}$ and $C_{[\lambda]}$ be the row and column stabilizers of $[\lambda]$, i.e.,
\begin{gather*}
R_{[\lambda]}= \{ f \in \sym{[\lambda]}\colon \forall\, (i, j) \in [\lambda],\, \exists\, (i,j')\in [\lambda],\, f(i,j) = (i,j') \}, \\
C_{[\lambda]}= \{ f \in \sym{[\lambda]}\colon \forall\, (i, j) \in [\lambda],\, \exists\, (i',j)\in [\lambda],\, f(i,j) = (i',j) \}.
\end{gather*}
Under the above group isomorphism, $R_{[\lambda]}$ corresponds to the row stabilizer $R_{\t}$ of $\t$, i.e., $R_{\t} = \t \circ R_{[\lambda]} \circ \t^{-1} \subseteq \sym{n}$, and similarly, $C_{[\lambda]}$ corresponds to the column stabilizer $C_{\t}$ of $\t$, i.e., $C_{\t}=\t\circ C_{[\lambda]}\circ \t^{-1}$.

Let $\t^{\lambda}$ denote the \textit{initial $\lambda$-tableau}, defined by $\t^{\lambda}(i,j) = \lambda_1 + \dotsb + \lambda_{i-1} + j$ for all $(i,j) \in [\lambda]$. Observe that $R_{\t^{\lambda}} = \sym{\lambda}$.

Post-composition of $\lambda$-tableaux by elements of the symmetric group $\sym{n}$ gives a well-defined, faithful and transitive left action of $\sym{n}$ on $\TT(\lambda)$, i.e., $\sigma \cdot \t = \sigma \circ \t$ for all $\sigma \in \sym{n}$ and $\t \in \TT(\lambda)$. Observe that \begin{gather*}R_{\sigma \cdot \t} = R_{\sigma \circ \t} = (\sigma \circ \t) \circ R_{[\lambda]} \circ (\sigma \circ \t)^{-1} = \sigma \circ R_{\t} \circ \sigma^{-1} = \sigma R_{\t} \sigma^{-1},\end{gather*} and similarly, $C_{\sigma \cdot \t} = \sigma C_{\t} \sigma^{-1}$. In particular, the row stabilizers of the $\lambda$-tableaux are conjugate subgroups of $\sym{\lambda}$.

A $\lambda$-tableau $\t$ is {\em standard} if it is increasing along each row and down each column, i.e., we have both $\t(i,j) < \t(i,j')$ and $\t(i,j)< \t(i',j)$ for all $(i,j),(i,j'),(i',j) \in [\lambda]$ with $j < j'$ and $i < i'$. Write $\TT_{\rm{std}}(\lambda)$ for the set of all standard $\lambda$-tableaux.

Let $(\alpha|\beta)$ be a bicomposition of $n$ and $\lambda$ be a partition of $n$. We view a $\lambda$-tableau $\T$ of type $(\alpha|\beta)$ as a colouring of the nodes of $[\lambda]$ by the colours $\bb{c}_1,\bb{c}_2,\dotsc, \bb{c}_{\ell(\alpha)},\bb{d}_1,\bb{d}_2,\dotsc,\bb{d}_{\ell(\beta)}$ such that there are exactly $\alpha_i$ nodes of colour $\bb{c}_i$ and $\beta_j$ nodes of colour $\bb{d}_j$ for all $1 \leq i\leq \ell(\alpha)$ and $1\leq j\leq \ell(\beta)$. Formally, $\T$ is a function \begin{gather*}\T\colon \ [\lambda]\to \{\bb{c}_1,\bb{c}_2,\dotsc, \bb{c}_{\ell(\alpha)},\bb{d}_1,\bb{d}_2,\dotsc,\bb{d}_{\ell(\beta)} \}\end{gather*} such that $|\T^{-1}(\{\bb{c}_i\})|=\alpha_i$ and $|\T^{-1}(\{\bb{d}_j\})|=\beta_j$ for all $1 \leq i\leq \ell(\alpha)$ and $1\leq j\leq \ell(\beta)$. We write $\TT(\lambda, (\alpha|\beta))$ for the set of all $\lambda$-tableaux of type $(\alpha|\beta)$.

Let $\t_0 \in \TT(\lambda)$ be a fixed $\lambda$-tableau (which may or may not be the initial $\lambda$-tableau $\t^{\lambda}$ defined above). We define the canonical $\lambda$-tableau of type $(\alpha|\beta)$ associated to $\t_0$, denoted as $\T_0$, as follows
\begin{gather*}
\T_0(i,j) = \begin{cases}\bb{c}_k&\text{if $\displaystyle{\sum^{k-1}_{a=1} \alpha_a < \t_0(i,j) \leq \sum^{k}_{a=1} \alpha_a}$ for some $k$,}\\ \bb{d}_k&\text{if $\displaystyle{|\alpha|+\sum^{k-1}_{b=1} \beta_b < \t_0(i,j) \leq |\alpha|+\sum^{k}_{b=1} \beta_b}$ for some $k$,}\end{cases} \end{gather*} namely, the nodes labelled $1,\ldots,\alpha_1$ by $\t_0$ are coloured $\bb{c}_1$, those labelled $\alpha_1 + 1,\dotsc, \alpha_1 + \alpha_2$ are coloured $\bb{c}_2$, and so on.

We also have a transitive left action on $\TT(\lambda, (\alpha|\beta))$ by $\sym{n}$ through $\t_0$ as follows: $\sigma \cdot \T = \T \circ (\t_0)^{-1} \circ \sigma^{-1} \circ \t_0$ for all $\sigma \in \sym{n}$ and $\T \in \TT(\lambda,(\alpha|\beta))$. Thus, if $(i,j) \in [\lambda]$ is labelled $a$ by $\t_0$, while the node labelled $\sigma^{-1}(a)$ by $\t_0$ is coloured $\bb{c}_k$ (respectively, $\bb{d}_k$) by $\T$, then $\sigma \cdot \T$ colours $(i,j)$ with $\bb{c}_k$ (respectively, $\bb{d}_k$). Observe that under this action, the stabiliser of $\T_0$ is $\sym{\alpha|\beta}$. Consequently, $\TT(\lambda, (\alpha|\beta))$ is in a bijective correspondence with the set of left cosets of $\sym{\alpha|\beta}$ in~$\sym{n}$ via $d \cdot \T_0 \leftrightarrow d \sym{\alpha|\beta}$. We set \begin{gather*}\T_d=d\cdot \T_0\end{gather*} so that $\T_d=\T_{d'}$ if and only if $d^{-1}d'\in\sym{\alpha|\beta}$.

We order the colours $\bb{c}_1,\bb{c}_2,\dotsc,\bb{d}_1,\bb{d}_2,\dotsc$ as follows
\begin{gather*}\bb{c}_1<\bb{c}_2<\cdots<\bb{d}_1<\bb{d}_2<\cdots.\end{gather*} A $\lambda$-tableau $\T$ of type $(\alpha|\beta)$ is {\em row semistandard} (respectively, {\em column semistandard}) if every row (respectively, column) of $\T$ is weakly increasing. More precisely, $\T$ is row semistandard if $\T(i,j) \leq \T(i, j')$ for all $(i,j), (i,j') \in [\lambda]$ with $j < j'$, and $\T$ is column semistandard if $\T(i,j) \leq \T(i', j)$ for all $(i,j), (i',j) \in [\lambda]$ with $i < i'$. Following \cite[Section~1.2]{Se} (see also \cite[Section~4]{DR}), we say that
a $\lambda$-tableau $\T$ of type $(\alpha|\beta)$ is {\em semistandard} if and only if
\begin{enumerate}\itemsep=0pt
\item [(i)] $\T$ is both row and column semistandard;
\item [(ii)] $\T(i,j) = \T(i,j')$ for some $(i,j),(i,j') \in [\lambda]$ with $j\ne j'$ only if $\T(i,j) = \bb{c}_k$ for some $k$;
\item [(iii)] $\T(i,j) = \T(i',j)$ for some $(i,j),(i',j) \in [\lambda]$ with $i\ne i'$ only if $\T(i,j) = \bb{d}_k$ for some $k$.
\end{enumerate}
We denote the set of semistandard $\lambda$-tableaux of type $(\alpha|\beta)$ by $\TTS(\lambda, (\alpha|\beta))$.

